\begin{trapbox}

	\begin{description}[style=nextline, leftmargin=2.8em, labelsep=0.5em, itemsep=0.35em]
		\item[\textbf{\textcolor{CTrap}{易错点一（定义域平移封闭）}}]
			有些同学只记住了周期公式$T=|a|$或$T=2|a|$，但忽略等式成立的前提：函数定义域必须对平移常数封闭。例如，$f(x)=\sqrt{x}$在$x=0$附近不一定满足周期性。
		\item[\textbf{\textcolor{CTrap}{正确理解（封闭性必需）：}}]
			周期函数的定义域$D$必须满足：若$x\in D$，则$x\pm T\in D$，其中$T$为周期。否则，即使给定等式，也无法定义周期。
		\item[\textbf{\textcolor{CTrap}{错因分析（忽视前提）：}}]
			机械记忆公式，忽视周期性定义中的定义域要求。
	\end{description}

	\medskip
	\begin{description}[style=nextline, leftmargin=2.8em, labelsep=0.5em, itemsep=0.35em]
		\item[\textbf{\textcolor{CTrap}{易错点二（周期长度误判）}}]
			对于$f(x+a)=-f(x)$，误以为周期是$|a|$，直接写$T=|a|$。
		\item[\textbf{\textcolor{CTrap}{正确理解（周期为2|a|）：}}]
			通过迭代可得$f(x+2a)=-f(x+a)$，结合$f(x+a)=-f(x)$得$f(x+2a)=f(x)$，周期应为$2|a|$。
		\item[\textbf{\textcolor{CTrap}{错因分析（漏乘因子2）：}}]
			只看到了$f(x+a)$与$f(x)$符号相反，认为$|a|$就能使函数值重复，忽略了符号变化后还需一次才能回到原值。
	\end{description}

	\medskip
	\begin{description}[style=nextline, leftmargin=2.8em, labelsep=0.5em, itemsep=0.35em]
		\item[\textbf{\textcolor{CTrap}{易错点三（符号识别错误）}}]
			在$f(x+a)=f(x+b)$中，对$a,b$的符号判断错误。例如由$f(x+1)=f(x-2)$，误得出$T=|1-2|$，即$T=1$。
		\item[\textbf{\textcolor{CTrap}{正确理解（a,b符号正确）：}}]
			公式$T=|a-b|$，需正确确定$a$和$b$。这里$a=1$，$b=-2$，得$T=|1-(-2)|$，即$T=3$。
		\item[\textbf{\textcolor{CTrap}{错因分析（符号误读）：}}]
			把$x-2$看作$x+(-2)$，容易忽视$b=-2$的负号，导致相减时出错。
	\end{description}

\end{trapbox}

\clearpage
\begin{trapbox}
\textbf{新增易错：得到 $2|a|$，就认定是最小正周期}
\begin{center}\begin{tikzpicture}[font=\small,>=stealth]
\draw[gray!40,rounded corners] (0,0) rectangle (14,6.8);
\node at (7,6.25) {$f(t)=\sin(3\pi t)$：平移 $1$ 变号，但 $2$ 不是最小正周期};
\begin{scope}[shift={(1.4,3.1)},x=4.9cm,y=1.3cm]
\draw[->,gray] (0,0)--(2.25,0) node[right]{$t$};
\draw[->,gray] (0,-1.2)--(0,1.8) node[above]{$y$};
\draw[blue,very thick,domain=0:2,samples=200] plot(\x,{sin(540*\x)});
% Horizontal distances are labeled by the two period arrows.
\draw[<->,orange!85!black] (0,1.45)--(.6666666667,1.45) node[midway,above]{$T_0=\tfrac23$};
\draw[<->,teal!80!black] (0,-1.5)--(2,-1.5) node[midway,below]{$T=2=3T_0$};
\node[below left=4pt] at (0,0){$0$};
\end{scope}
\node at (7,.45) {连续代入能证明“一个周期”；最小正周期还要另作判断};
\end{tikzpicture}
\end{center}
图中 $f(x+1)=-f(x)$，所以 $2$ 是一个周期；但最小正周期是 $2/3$。
\textbf{正确写法：}连续代入得到的 $2|a|$ 是一个正周期；只有排除了全部更小的正周期，才能称为最小正周期。

\medskip
\textbf{新增易错：分式只检查第一次分母}

第一式需要 $f(x)-1\ne0$，第二式需要 $1+f(x)\ne0$；等式必须在全部定义域内成立，且定义域允许连续平移。不能只在某一点算出“两次还原”，就断言整个函数有周期。两篇新增证明已分别核验第二次分母仍然非零。
\end{trapbox}
